Un rayo de luz viaja en línea recta desde el punto $O = (0, 0, 0)$ en la dirección $\vec{d} = (1, 2, 2)$. Una varilla delgada está representada por el segmento que une los puntos $P = (3, 1, 0)$ y $Q = (3, 5, 4)$.
\begin{enumerate}
    \item Parametrizar la trayectoria del rayo y el segmento de la varilla.
    \item Determinar si el rayo intersecta el segmento de la varilla. Si es así, encontrar el punto exacto de intersección.
    \item Calcular la distancia del origen $O$ a la recta que contiene la varilla.
    \item Normalizar el vector director $\vec{d}$ del rayo y usar los cosenos directores para describir el ángulo que forma con cada eje coordenado.
\end{enumerate}